\chapter{استنتاج‌های ریاضی}
\label{app:math}

این پیوست استنتاج‌های ریاضی برای اجزای کلیدی را ارائه می‌کند.

%==============================================================================
\section{تعریف \lr{AUROC}}
%==============================================================================

مساحت زیر منحنی \lr{ROC} (\lr{AUROC}) به‌صورت زیر تعریف می‌شود:
\begin{equation}
\text{\lr{AUROC}} = \int_0^1 \text{\lr{TPR}}(t) \, d(\text{\lr{FPR}}(t)) = \int_0^1 \text{\lr{TPR}}(\text{\lr{FPR}}^{-1}(u)) \, du
\end{equation}
که در آن \lr{TPR} (نرخ مثبت صحیح) و \lr{FPR} (نرخ مثبت کاذب) توابعی از آستانه‌ی تصمیم $t$ هستند. به‌طور معادل، \lr{AUROC} برابر است با احتمال اینکه یک نمونه‌ی مثبت تصادفی بالاتر از یک نمونه‌ی منفی تصادفی رتبه‌بندی شود:
\begin{equation}
\text{\lr{AUROC}} = P(\text{\lr{score}}(X^+) > \text{\lr{score}}(X^-))
\end{equation}
\lr{AUROC} نسبت به آستانه ثابت است و برای طبقه‌بندی نامتعادل مناسب است.

%==============================================================================
\section{امتیاز \lr{F1} کلان}
%==============================================================================

برای طبقه‌بندی دودویی با کلاس‌های \pnum{0} و \pnum{1}:
\begin{equation}
F1_k = \frac{2 \cdot \text{\lr{Precision}}_k \cdot \text{\lr{Recall}}_k}{\text{\lr{Precision}}_k + \text{\lr{Recall}}_k}
\end{equation}
امتیاز \lr{F1} کلان میانگین روی کلاس‌هاست:
\begin{equation}
F1_{\text{\lr{macro}}} = \frac{F1_0 + F1_1}{2}
\end{equation}
این هر دو کلاس را به‌طور یکسان در نظر می‌گیرد؛ برای داده‌ی نامتعادل مهم است.

%==============================================================================
\section{ضریب همبستگی متیوز (\lr{MCC})}
%==============================================================================

\lr{MCC} به‌صورت زیر تعریف می‌شود:
\begin{equation}
\text{\lr{MCC}} = \frac{TP \cdot TN - FP \cdot FN}{\sqrt{(TP+FP)(TP+FN)(TN+FP)(TN+FN)}}
\end{equation}
\lr{MCC} در بازه‌ی $[-1, 1]$ است؛ مقدار \pnum{0} پیش‌بینی تصادفی را نشان می‌دهد. متقارن است و نسبت به نامتعادل کلاس‌ها مقاوم است.

%==============================================================================
\section{گرادیان تابع زیان شبکه‌ی نمونه‌اولیه‌ای (\lr{Prototypical})}
%==============================================================================

برای شبکه‌های نمونه‌اولیه‌ای با تابع زیان آنتروپی متقاطع $\mathcal{L} = -\log p(y|x)$، گرادیان نسبت به بردار نهفتۀ پرس‌وجو $f_\theta(x)$ عبارت است از:
\begin{equation}
\frac{\partial \mathcal{L}}{\partial f_\theta(x)} = \frac{1}{\tau} \left( \mathbf{c}_{y} - \sum_{k} p(y=k|x) \mathbf{c}_k \right)
\end{equation}
که در آن $\mathbf{c}_k$ نمونه‌اولیه‌ی کلاس $k$ است. گرادیان پرس‌وجو را به سمت نمونه‌اولیه‌ی واقعی می‌کشد و از بقیه دور می‌کند.

%==============================================================================
\section{زیان کانونی}
%==============================================================================

زیان کانونی  نمونه‌های آسان را کم‌اهمیت می‌کند:
\begin{equation}
\mathcal{L}_{\text{\lr{focal}}} = -\alpha (1-p_t)^\gamma \log(p_t)
\end{equation}
که در آن $p_t = p(y|x)$ احتمال پیش‌بینی‌شده برای کلاس درست است. برای $\gamma > 0$، نمونه‌های آسان ($p_t \approx 1$) سهم کمتری دارند. ما $\gamma = 2$ استفاده می‌کنیم.

%==============================================================================
\section{قانون به‌روزرسانی \lr{FedAvg}}
%==============================================================================

میانگین‌گیری فدرال (\lr{Federated Averaging}) به‌روزرسانی‌های محلی را با میانگین وزنی تجمیع می‌کند:
\begin{equation}
\theta_{t+1} = \sum_{k=1}^{K} \frac{n_k}{n} \theta_{t+1}^{(k)}
\end{equation}
که در آن $n_k$ تعداد نمونه‌ها در کلاینت $k$، $n = \sum_k n_k$، و $\theta_{t+1}^{(k)}$ مدل پس از آموزش محلی روی کلاینت $k$ است. وقتی داده‌ها \lr{IID} باشند، گرادیان مورد انتظار هدف سراسری حفظ می‌شود؛ تحت \lr{non-IID} همگرایی ممکن است کندتر باشد.

%==============================================================================
\section{خودتوجهی}
%==============================================================================

برای دنباله‌ی ورودی $\mathbf{X} \in \mathbb{R}^{L \times d}$، خودتوجه به‌صورت زیر محاسبه می‌شود:
\begin{equation}
\text{\lr{Attention}}(\mathbf{Q}, \mathbf{K}, \mathbf{V}) = \text{\lr{softmax}}\left(\frac{\mathbf{Q}\mathbf{K}^\top}{\sqrt{d_k}}\right) \mathbf{V}
\end{equation}
که در آن $\mathbf{Q} = \mathbf{X}\mathbf{W}_Q$، $\mathbf{K} = \mathbf{X}\mathbf{W}_K$، $\mathbf{V} = \mathbf{X}\mathbf{W}_V$. توجه چندسر $h$ چنین عملیاتی را موازی اجرا کرده و خروجی‌ها را به‌هم می‌چسباند. در \lr{ETHOS} از $h=8$ سر استفاده می‌کنیم.

%==============================================================================
\section{بازه‌ی اطمینان \lr{bootstrap}}
%==============================================================================

برای \lr{AUROC}، بازه‌ی اطمینان \pnum{95}\% با نمونه‌برداری مجدد از مجموعه‌ی آزمون $B=1000$ بار با جایگزینی محاسبه می‌کنیم. برای هر نمونه‌ی \lr{bootstrap} $b$، $\hat{A}_b = \text{\lr{AUROC}}(y, \hat{y}_b)$ را می‌سنجیم. بازه‌ی اطمینان $[\hat{A}_{(25)}, \hat{A}_{(975)}]$ است که زیرنویس‌ها صدک‌های توزیع \lr{bootstrap} را نشان می‌دهند.

%==============================================================================
\section{آزمون \lr{DeLong}}
%==============================================================================

آزمون \lr{DeLong}~\cite{delong1988comparing} دو منحنی \lr{ROC} همبسته را مقایسه می‌کند. برای پیش‌بینی‌های جفت $(s_1, s_2)$ روی همان نمونه‌ها، آماره‌ی آزمون بر اساس تفاوت در \lr{AUC} تجربی است. تحت فرض صفر (برابری \lr{AUC})، آماره به‌طور مجانبی نرمال است. در این پایان‌نامه فقط در مواردی می‌توان چنین آزمونی را تفسیر کرد که دو مدل روی همان نمونه‌های آزمون و با بردارهای احتمال هم‌تراز ارزیابی شده باشند؛ بنابراین مقایسه‌های بین‌پروتکلی به‌صورت توصیفی گزارش می‌شوند و آزمون‌های جفتی کامل به کار آینده سپرده می‌شوند.
